\documentclass[11pt,a4paper]{article}
\usepackage[T1]{fontenc}
\usepackage[utf8]{inputenc}
\usepackage[spanish,es-noshorthands]{babel}
\usepackage[margin=2.5cm]{geometry}
\usepackage{enumitem}
\setlist{itemsep=2pt}

\title{Pedestal y ruido del Citiroc con la CAEN A5202}
\author{}
\date{\today}

\begin{document}
\maketitle

\noindent La tarjeta que usamos es una CAEN A5202 (dos Citiroc~1A, 64 canales),
controlada con el software Janus. Los nombres de los parámetros son los de
Janus (Sección~\ref{sec:janus}). Estudio número 1.

\section*{Programa de festejos}

\noindent {\bf Setup}
\begin{itemize}
\item {\bf Acción:} montar sin cristal, con los  SiPM conectados y a oscuras, polarizados entre 1 y 2~V por
  debajo de la tensión de ruptura (breakdown voltage, BV;
  Sección~\ref{sec:razon}). En Janus,
  usar los parámetros definidos en la Sección~\ref{sec:janus}.
\item {\bf Resultado:} cada run se para solo a los 10\,000 eventos y lleva su
  propio número \texttt{N}. Al acabar, Janus escribe en \texttt{DataFilePath}
  \texttt{RunN\_list.csv} (para cada evento y canal, la baja y la alta
  ganancia), los histogramas de cada canal \texttt{RunN\_PHA\_HG\_0\_ch.txt}
  y \texttt{RunN\_PHA\_LG\_0\_ch.txt}, \texttt{RunN\_Info.txt} (la
  configuración del run) y el fichero de temperatura y tensión.
\item {\bf Precaución:} en el primer run, comprobar que los SiPM no dan avalanchas: el espectro de
  alta ganancia de cada canal debe ser una sola gaussiana, sin pico de
  1 fotoelectrón (pe) a su derecha. Si aparece, bajar la tensión 0,5~V y
  repetir.
\end{itemize}

\noindent {\bf Paso 1: ruido frente al shaping de la alta ganancia}
\begin{itemize}
\item {\bf Acción:} un run por cada valor de \texttt{HG\_ShapingTime} =
  25, 50 y 87,5~ns, con \texttt{LG\_ShapingTime} = 87,5~ns en todos:
  3 runs.
\item {\bf Resultado:} 3 ficheros \texttt{RunN\_list.csv}, con sus
  histogramas.
\item {\bf Análisis:} el de la Sección~\ref{sec:datos}, para cada run
  y canal, la media (pedestal) y la anchura (ruido) de la alta y la baja
  ganancia.
\end{itemize}

\noindent {\bf Paso 2 (opcional): pedestal en peak detector}
\begin{itemize}
\item {\bf Acción:} un run con \texttt{PeakDetectorMode} = PEAK\_STRETCH,
  \texttt{HG\_ShapingTime} = 25~ns, \texttt{LG\_ShapingTime} = 87,5~ns y
  \texttt{HoldDelay} = 200~ns.
\item {\bf Resultado:} el pedestal tal como lo verán los runs de energía
  (los runs con cristal y fuente), que usan peak detector: más alto y más estrecho que en track \& hold,
  porque el chip guarda el máximo del ruido dentro de la ventana.
\end{itemize}

\noindent {\bf Qué necesita tu estudiante}
\begin{itemize}
\item Todo lo que hay en \texttt{DataFilePath}.
\item Logbook, una línea por run: número de run, paso al que
  corresponde y cualquier incidencia. La configuración, la temperatura y
  la tensión ya están en los ficheros de Janus.
\end{itemize}

\section{Cómo entra la medida en el modelo}
\label{sec:razon}

\begin{itemize}
\item El ruido electrónico de cada canal entra en la resolución en energía y
  en la anchura de los picos de los espectros de un fotoelectrón (single
  photoelectron, SPE) con los que se calibra la ganancia. La simulación
  estima el ruido, pero la estimación puede quedarse corta. Necesitamos datos. 
\item El ruido de un preamplificador crece con la capacitancia de su
  entrada, y la de nuestros SiPM es grande, unos 2000~pF por canal. Por eso hay que medirlo con los SiPM conectados.
\item Por debajo del BV el SiPM conserva su capacitancia
  pero no produce avalanchas. El pedestal medido así
  contiene, por tanto, solo el ruido electrónico. Por encima del BV aparecen también
  los pulsos oscuros (dark pulses, DP), con una frecuencia de varios MHz por canal.
\item El BV es 38,0~V a 25~$^\circ$C y baja 34~mV por
  grado; a temperatura ambiente, \texttt{HV\_Vbias} $\approx$ 36,5~V deja todos
  los canales entre 1 y 2~V por debajo. 
  \item El modelo de la simulación supone 0,1~pe de ruido en alta ganancia
  con shaping de 87,5~ns y 3~V de overvoltage (tensión por encima del BV).
  El ruido es fijo en carga, y la carga de un pe es proporcional al
  overvoltage, así que a 5~V son 0,06~pe. Supone además que el ruido crece como $(87{,}5~\mathrm{ns}/t_p)^{1/2}$
  al acortar el peaking time $t_p$, y que la baja ganancia tiene el mismo
  ruido en cuentas de ADC que la alta. Estas medidas nos permiten entender lo realistas que son estas suposiciones.
\end{itemize}

\section{Parámetros de Janus (A5202)}
\label{sec:janus}

\begin{center}
\small
\begin{tabular}{p{0.3\textwidth} p{0.6\textwidth}}
\hline
parámetro & valor \\
\hline
\texttt{HV\_Vbias} & el que deje la tensión en cada SiPM entre 1 y 2~V por
debajo del BV (unos 36,5~V a temperatura ambiente si \texttt{HV\_IndivAdj}
no resta nada); \texttt{HV\_IndivAdj}: los valores que tenga cargados
la tarjeta, los mismos en los tres estudios \\
\texttt{AcquisitionMode} & SPECTROSCOPY \\
\texttt{BunchTrgSource} & PTRG (trigger periódico, independiente de la señal) \\
\texttt{PtrgPeriod} & 1~ms: 10\,000 eventos en 10~s \\
\texttt{PeakDetectorMode} & TRACK\&HOLD (paso 1); PEAK\_STRETCH (paso 2) \\
\texttt{HG\_ShapingTime} & 25, 50 o 87,5~ns, según el run \\
\texttt{LG\_ShapingTime} & 87,5~ns \\
\texttt{HoldDelay} & cualquiera en el paso 1 (sin señal no importa); 200~ns en el paso 2 \\
\texttt{HG\_Gain} & el valor por defecto (50) \\
\texttt{LG\_Gain} & el valor por defecto (50) \\
\texttt{GainSelect} & BOTH \\
\texttt{StopRunMode}, \texttt{PresetCounts} & PRESET\_COUNTS, 10000 \\
\texttt{OF\_ListCSV}, \texttt{OF\_SpectHisto}, \texttt{OF\_RunInfo},
\texttt{OF\_ServiceInfo} & 1 \\
\texttt{RunNumber\_AutoIncr} & 1 \\
\hline
\end{tabular}
\end{center}

\section{Análisis}
\label{sec:datos}

\begin{itemize}
\item Para cada canal y cada ganancia, una gaussiana
  ajustada al espectro. Su media es el pedestal y su sigma el ruido, los
  dos en cuentas de ADC. Para comparar runs se usa la mediana sobre los 64
  canales. Un canal con una cola o un segundo pico tiene avalanchas o
  ruido de otro origen, y se anota aparte.
\item {\bf Paso a fotoelectrones:} el ruido en pe es el ruido en cuentas de
  ADC dividido por la ganancia en cuentas de ADC por pe, que miden los
  runs de calibración con SPE. Hay que usar la ganancia medida con el mismo
  shaping y el mismo \texttt{HG\_Gain}.
\end{itemize}

\section{Lo que se espera, según el modelo}

\begin{itemize}
\item {\bf Ruido de la alta ganancia, en pe de 5~V} (los de la calibración):
  0,06 a 87,5~ns, 0,08 a 50~ns y 0,11 a 25~ns.
\item {\bf Prueba sin ganancia:} el cociente entre el ruido a un shaping y
  el ruido a 87,5~ns, en cuentas de ADC, no depende de la ganancia. El
  modelo da 1,32 a 50~ns y 1,87 a 25~ns.
\item {\bf Baja ganancia:} el mismo ruido en cuentas de ADC que la alta a
  87,5~ns (run de 87,5~ns del paso 1, donde las dos tienen el mismo
  shaping).
\item {\bf Pedestal:} distinto en cada canal, estable entre runs. En peak
  detector (paso 2), más alto que en track \& hold.
\end{itemize}

\end{document}
